% META: {"id":"1SPE-TRIGO-CO-007","chapitre":"1SPE-TRIGONOMETRIE","type_objet":"corrige","exercice_ref":"1SPE-TRIGO-EX-007","capacites_codes":["C1"],"sources_inspiration":[],"status":"approved","origin":"generated","status_history":["generated","audited","accepted","approved"],"release_acceptance":"RELEASE_OWNER_BATCH_ACCEPTANCE","acceptance_closure_digest":"sha256:9b3ccf9a81c5520fbb7e03b7b2d3e7bf2057a02834b6d908bf3be5f49f3b553f","acceptance_receipt_digest":"sha256:4fad86f0282e0fa4e0419cca1ad6379ae7af5a280c04a4a90880f90a1c044242"}
% BEGIN-VERIFY
% from sympy import *
% R = 6371
% nm = R * pi / (60 * 180)
% assert abs(float(nm) - 1.853) < 0.01
% assert 5 * 60 == 300
% END-VERIFY
\begin{corrige}{1SPE-TRIGO-EX-007}

\textbf{1.} $1' = \dfrac{1}{60}° = \dfrac{1}{60} \times \dfrac{\pi}{180} = \dfrac{\pi}{10800}$ rad.

\textbf{2.} La longueur d'un mille nautique est $\ell = R\theta = 6371 \times \dfrac{\pi}{10800} \approx 1{,}852$~km.

\textbf{3.} $5° = 5 \times 60 = 300$ minutes d'arc, soit $300$ milles nautiques.

\end{corrige}
